% Generated draft: requires mathematical and source review.
\documentclass[11pt]{article}
\usepackage[margin=25mm]{geometry}
\usepackage{fontspec}
\setmainfont{Latin Modern Roman}
\usepackage{amsmath,amssymb,mathtools,unicode-math}
\setmathfont{Latin Modern Math}
\usepackage{enumitem,needspace,xurl,hyperref}
\setlength{\parindent}{0pt}
\setlength{\parskip}{5pt}
\setlength{\emergencystretch}{3em}
\newcommand{\N}{\mathbb{N}}
\newcommand{\Z}{\mathbb{Z}}
\newcommand{\Q}{\mathbb{Q}}
\newcommand{\R}{\mathbb{R}}
\newcommand{\C}{\mathbb{C}}
\newcommand{\F}{\mathbb{F}}
\newcommand{\A}{\mathbb{A}}
\newcommand{\PP}{\mathbb{P}}
\newcommand{\E}{\mathbb{E}}
\newcommand{\eps}{\varepsilon}
\newcommand{\dd}{\,\mathrm{d}}
\begin{document}
\section*{UnsolvedMath ID: 20000441}
\textbf{Source problem code:} catalogue entry
\subsection{Definitions and mathematical statement}
The source asks to use intermediate Jacobians arithmetically and, in dimension three, to study algebraic one-cycles modulo rational equivalence. Its concrete target is a smooth projective rigid Calabi--Yau threefold \(X/\C\), with
\[
K_X\sim0,\quad\pi_1(X)=1,\quad
h^{1,0}=h^{2,0}=h^{2,1}=0,\quad h^{3,0}=1.
\]
Determine the elliptic curve
\[
J^2(X)=H^3(X,\C)/(F^2H^3(X,\C)+H^3(X,\Z)_{\mathrm{free}}),
\]
in particular its \(j\)-invariant. Here \(F^2=H^{3,0}\), so this complex torus has dimension one. The source does not specify an arithmetic field of definition or an effective period/lattice input, and determining this torus does not compute all of \(\mathrm{CH}^2(X)\).

The attached aid reports the exact period expression. For a primitive integral basis \(\alpha,\beta\) of torsion-free \(H^3\), \(0\ne\Omega\in H^{3,0}\), and cup-product pairing \(Q\), choose orientation with
\[
\tau=\frac{Q(\beta,\Omega)}{Q(\alpha,\Omega)}\in\mathbb H,
\qquad j(J^2(X))=j(\tau).
\]
A rational basis or isogeny comparison alone need not determine this exact elliptic curve. The aid also reports \(J_a^3(X)=0\), where this notation means the connected algebraic Abel--Jacobi part generated by algebraically trivial one-cycles. This does not assert the full Abel--Jacobi map on homologically trivial cycles or the Chow group is zero, and does not give an automatic arithmetic descent of all of \(J^2(X)\).

For a genuine lattice-preserving automorphism acting on \(H^{3,0}\) by a primitive third or sixth root, the aid reports \(j=0\); a primitive fourth root gives \(j=1728\). A primitive fifth-root scalar action is impossible on rational rank-two \(H^3\). It consequently reports \(j=0\) for the standard 27-fixed-point CM quotient example. The source's description using an elliptic curve with CM by \(\Q(\zeta_3)\) and kernels of \(\zeta_3-1\) omits the quotient action: the product of finite kernels alone is not a threefold. The usual quotient/resolution interpretation is therefore explicit but not fully specified by the extracted text.

The source also gives the nodal quintic
\[
\sum_{i=1}^5x_i^5-5\psi\prod_{i=1}^5x_i=0,\qquad\psi^5=1,
\]
with 125 nodes and a rigid resolution, attributed to Yui. The aid calls it the Schoen quintic and says its visible fifth-root symmetry does not determine \(j\); a primitive integral lattice comparison is still needed. No \(j\)-value for it is supplied.

Its other remark concerns smooth intersections of two quadrics in \(\PP^5\): study the genus-two curve \(y^2=\det(\lambda f_1+\mu f_2)\), its Jacobian, and the variety of lines as an abelian-variety torsor. It reports a line over finite fields from that torsor description, and asks about arithmetic applications, skew conjugate lines, ground-field conics and finding a line over a quadratic extension. The determinant model and finite-field claims need their smoothness/characteristic hypotheses; they are separate from the rigid Calabi--Yau computation.
\subsection{Short English statement}
Determine the elliptic intermediate Jacobian of a rigid Calabi--Yau threefold, ideally through its exact \(j\)-invariant, and explore arithmetic uses. The attached aid gives a primitive-lattice period formula and limited symmetry certificates, distinguishes the zero algebraic Abel--Jacobi part from the full torus, and leaves the quintic's exact lattice comparison unresolved.
\subsection{Sources}
\begin{itemize}
\item Local UnsolvedMath registry, permanent ID \(20000441\), exact public code \texttt{catalogue entry}; canonical location \texttt{aim-algebraic-number-theory-notes.json}, notes[92].
\item AIM workshop ``Rational and integral points on higher dimensional varieties,'' Question 9, source version dated 22 November 2004; supplied URL \url{https://aimath.org/WWN/qptsurface2/qptsurface2.pdf}. Remarks credit Colliot-Th\'el\`ene, ``Candela'' as extracted, and Yui. An incomplete Harari web-footnote fragment, with domain \texttt{www.dma.ens.fr}, is retained in the excerpt. The PDF was not consulted here; the literature field is empty.
\item Attached research attempt 1, ``Rigid Calabi--Yau intermediate Jacobians: periods, cycle-theoretic separation, and symmetry certificates,'' without full bibliographic citations for its symmetry/quotient claims or Schoen-quintic lattice problem.
\end{itemize}
\end{document}
